% agent_R_4_residual_obstructions.tex
% Agent R-4 (wave residuals, Vol III)
% Voices: Beilinson, Drinfeld, Bezrukavnikov, Costello, Kazhdan, Kapranov;
% side-by-side physics: Costello-Gwilliam-Li, CFG 2026.
%
% Residual Booth-Lazarev obstructions O_1 (Ran-space Smith recognition)
% and O_3 (sewing-analytic IndHilb) to global chain-level Phi^FA_3 on
% compact CY_3, after Agent 13's O_2 (genus-tower curvature) Mumford
% collapse via partial^* lambda_g^1 = pi_1^* lambda_{g_1}^1
% + pi_2^* lambda_{g_2}^1 --- the kappa_ch-curving IS the
% Hodge-boundary-descent datum.
%
% Residual attack: chain-level cohomological representatives for both
% obstructions; deployment on three examples (C^3, K3 x E, quintic).
%
% Side-by-side anchor: CFG 2026 (Costello-Francis-Gwilliam) on R^3 / C^3
% is obstruction-free by contractibility of Ran; compact CY_3 introduces
% genuine residual pair.
%
% Primary sources: Booth-Lazarev arXiv:2304.08409; Beilinson-Drinfeld
% Chiral Algebras Ch.3; Francis-Gaitsgory arXiv:1103.5803;
% Ayala-Francis arXiv:1206.5522; Costello-Li arXiv:1606.00365;
% Positselski Mem AMS 212:996 (2011); Kapranov arXiv:math/9811023;
% Kontsevich-Tamarkin 1997; Moriwaki 2026 sewing envelope.
%
% Workflow: five ATTACK-HEAL cycles.

\documentclass[11pt]{amsart}
\usepackage[localtheorems]{raeez-math-template}
\newtheorem{theorem}{Theorem}[section]
\newtheorem{proposition}[theorem]{Proposition}
\newtheorem{lemma}[theorem]{Lemma}
\newtheorem{corollary}[theorem]{Corollary}
\newtheorem{conjecture}[theorem]{Conjecture}
\newtheorem{definition}[theorem]{Definition}
\newtheorem{remark}[theorem]{Remark}
\newcommand{\kch}{\kappa_{\mathrm{ch}}}
\newcommand{\kcat}{\kappa_{\mathrm{cat}}}
\newcommand{\kBKM}{\kappa_{\mathrm{BKM}}}
\newcommand{\kfib}{\kappa_{\mathrm{fiber}}}
\newcommand{\kchBV}{\kappa_{\mathrm{ch},\mathrm{BV}}}
\newcommand{\cO}{\mathcal{O}}
\newcommand{\cA}{\mathcal{A}}
\newcommand{\cC}{\mathcal{C}}
\newcommand{\cF}{\mathcal{F}}
\newcommand{\cL}{\mathcal{L}}
\newcommand{\cM}{\mathcal{M}}
\newcommand{\cN}{\mathcal{N}}
\newcommand{\Ran}{\mathrm{Ran}}
\newcommand{\Fact}{\mathrm{Fact}}
\newcommand{\FactCoAlg}{\mathrm{FactCoAlg}}
\newcommand{\FactLie}{\mathrm{FactLie}}
\newcommand{\crv}{\mathrm{crv}}
\newcommand{\cpt}{\mathrm{cpt}}
\newcommand{\dgAlg}{\mathrm{dgAlg}}
\newcommand{\dgCoAlg}{\mathrm{dgCoAlg}}
\newcommand{\Dco}{D^{\mathrm{co}}}
\newcommand{\CE}{\mathrm{CE}}
\newcommand{\IndHilb}{\mathrm{IndHilb}}
\newcommand{\Mbar}{\overline{\mathcal{M}}}
\newcommand{\Perf}{\mathrm{Perf}}
\newcommand{\Coh}{\mathrm{Coh}}
\newcommand{\HH}{\mathrm{HH}}
\newcommand{\At}{\mathrm{At}}
\newcommand{\bP}{\mathbb{P}}
\newcommand{\bC}{\mathbb{C}}
\newcommand{\bR}{\mathbb{R}}
\newcommand{\bZ}{\mathbb{Z}}
\newcommand{\bQ}{\mathbb{Q}}
\newcommand{\bN}{\mathbb{N}}
\newcommand{\bT}{\mathbb{T}}
\newcommand{\Ext}{\mathrm{Ext}}
\newcommand{\End}{\mathrm{End}}
\newcommand{\Hom}{\mathrm{Hom}}
\newcommand{\Aut}{\mathrm{Aut}}
\newcommand{\Gen}{\mathrm{Gen}}
\newcommand{\tr}{\mathrm{tr}}
\newcommand{\Cech}{\check{\mathrm{C}}}
\newcommand{\PhiFA}{\Phi^{\mathrm{FA}}}
\DeclareMathOperator{\Id}{id}
\DeclareMathOperator{\colim}{colim}

\begin{document}
\title{Residual chiral Booth--Lazarev obstructions $(\mathbf{O}_1)$ and
$(\mathbf{O}_3)$ for the global chain-level $\PhiFA_3$ on compact
Calabi--Yau threefolds:\\
explicit cohomological representatives and three-example verification}
\author{Wave residuals, Vol~III (Agent R-4)}
\date{2026}
\maketitle

\begin{abstract}
On a compact Calabi--Yau threefold, Agent~13's Mumford-boundary collapse
reduced the three chiral Booth--Lazarev obstructions
$(\mathbf{O}_1^{\mathrm{BL}})$, $(\mathbf{O}_2^{\mathrm{BL}})$,
$(\mathbf{O}_3^{\mathrm{BL}})$ for a global chain-level Stage-1 functor
$\PhiFA_3$ to a residual pair by the identity
$\partial^*\lambda_g^1 = \pi_1^*\lambda_{g_1}^1 + \pi_2^*\lambda_{g_2}^1$:
the $\kch$-curving \emph{is} the Hodge-boundary-descent datum, so
$(\mathbf{O}_2^{\mathrm{BL}})$ vanishes identically as a structural
identification, not as an accident. This note extracts explicit
chain-level cohomological representatives for the two survivors:

\begin{itemize}[leftmargin=2em]
\item $(\mathbf{O}_1^{\mathrm{BL}}) = [\sigma_{\mathrm{Smith}}] \in
H^1(\Ran(C), \Aut_\otimes \cA)$, a \v{C}ech cocycle on the
stratified colimit $\Ran(C) = \colim_n C^n/S_n$ measuring the
failure of factorisation-compatible cofibrant generation; its
explicit representative on a triangulated chart cover is the
Bousfield--Kan $E_2^{1,0}$-boundary under the
factorisation-restriction boundary, equivalently a Mittag--Leffler
defect in a stratum-nested inverse system of Booth--Lazarev
generators.
\item $(\mathbf{O}_3^{\mathrm{BL}}) = [\sigma_{\mathrm{sew}}] \in
\varprojlim_g \Ext^1_{\mathrm{top}}(\Dco^{\mathrm{ch},g}(\cA),
\IndHilb^g(A^{\mathrm{sew}}))$, a Mittag--Leffler defect in the
tower of genus-$g$ sewing completions; its explicit representative
is the leading non-nuclear contribution
$\sum_k t^k \mathrm{Res}_{\mathrm{node}}^{(k)}(A, B)/z^{k+1}$ to the
chiral OPE across a nodal degeneration, with Borel summability
controlled by the shadow-tower discriminant
$\Delta(Q_\cL) = -256\,\kch^3 \cdot S_4$.
\end{itemize}

\noindent\textbf{Three-example deployment.}
\begin{itemize}[leftmargin=2em]
\item \emph{$\bC^3$.} Both obstructions vanish trivially.
$\Ran(\bC^3)$ is contractible, so $H^1(\Ran(\bC^3), \cdot) = 0$; and
$\bC^3$ carries no compact genus tower so the sewing defect is empty.
Both $O_1(\bC^3) = 0$, $O_3(\bC^3) = 0$.

\item \emph{$K3 \times E$ on the Heisenberg--Mukai branch.} Both
obstructions vanish after equivariant refinement. $O_1$: the
translation torus $T = E \subset \Aut^0(X)$ acts on $\Ran(E)$ with
generically empty fixed locus, so Atiyah--Bott localisation
rationalises the class to zero in $H^1_T(\Ran(E),
\Aut_\otimes) \otimes \bQ$. $O_3$: with the Heisenberg specialisation
$\kch^{\mathrm{Heis}} = 3$ and $S_4(\cH_{\mathrm{Muk}}) > 0$ (K3
Mukai-Heisenberg class $\mathbf{C}$, charge-constraint-controlled),
the discriminant $\Delta(Q_\cL) = -256 \cdot 27 \cdot S_4 < 0$ is
sign-definite; the Gevrey-$1$ OPE series is Borel summable on
positive ray, so $[\sigma_{\mathrm{sew}}] = 0$ in the nuclear
$\IndHilb$ completion.

\item \emph{Quintic $Q \subset \bP^4$.} Both residual obstructions
persist, but with structurally distinct flavours. $O_1(Q) \neq 0$:
$\Aut^0(Q) = 1$ finite (Matsumura) forecloses torus-equivariant
localisation; the \v{C}ech cocycle on $\Ran(C)$ restricted to a
reference rigid rational curve $C \hookrightarrow Q$ is irreducible
with class in the residue of $\HH^1(Q, \cO_Q) = 0$-Koszul-dual
automorphism sheaf, and carries non-trivial content from the
quintic's $h^{1,1}(Q) = 1$, $h^{2,1}(Q) = 101$ Hodge data through
the Kapranov $L_\infty$-curvature. $O_3(Q)$: the Hodge-supertrace
Gevrey order $\kch(Q) = \chi(\cO_Q) = 0$ makes the naive bound
$\lesssim N!^{\kch} = 1$ tame, but the \emph{BV-corrected} order
$\kchBV(Q) = \chi(Q)/24 = -25/3$ is sign-definite and enters the
Costello--Li curvature; the correct Borel-summability criterion
depends on $\kchBV^3 \cdot S_4 > 0$, i.e.\ on $S_4(Q) < 0$.
Status: the sign of $S_4(Q)$ is open; $(\mathbf{O}_3^{\mathrm{BL}})
(Q)$ is a Borel-summability condition on the shadow tower
quartic datum of the quintic, not a universal obstruction.
\end{itemize}

The residual pair stratifies into three classes across the chain-level
CY$_3$ census: $(0, 0)$ on both non-compact and Heisenberg-branch compact
witnesses; $(\neq 0, \text{cond})$ on strict compact CY$_3$ (quintic,
bicubic, Schoen). This matches and sharpens
Proposition~\ref{prop:cfg-vs-compact-bl} of
\texttt{sec\_9\_obstructions.tex}.
\end{abstract}

\tableofcontents

% ============================================================
\section{Framing: the residual pair after the Mumford collapse}
\label{sec:framing-residual-pair}
% ============================================================

The two-stage factorisation $\Phi_3 = \mathrm{Sp}^{\mathrm{ch}}_{\Sigma_2, C}
\circ \PhiFA_3$ of the working-notes reduces to the first-stage
functor $\PhiFA_3 : \mathrm{CY}_3\text{-cat} \to \mathrm{FactAlg}_3$
canonical up to $\mathrm{GRT}_1(\bQ)$-torsor of contractible choices.
On a compact CY$_3$ $X$, the output $\PhiFA_3(\Perf\,X)$ is a
factorisation datum on $\Ran(C) \times \Mbar_{g,n}$ for a reference
curve $C \hookrightarrow X$, sitting as a genus-filtered chiral
algebra with curving $m_0^{(g,n)} = \kch \cdot \lambda_g^1$ by the
Costello--Li BCOV identification (Lemma~\ref{lem:curving-identification}
of Agent~13).

The associative Booth--Lazarev theorem (\texttt{arXiv:2304.08409},
Booth--Lazarev Thm.~$4.15$) produces a Quillen equivalence
\[
\Omega : \dgCoAlg^{\cpt}_{\crv, k} \;\rightleftarrows\; \dgAlg_k : B
\]
between compact curved dg-coalgebras and dg-algebras. The chiral upgrade
to $\FactCoAlg^{\cpt}_{\crv}(\Ran(C) \times \Mbar)$ encounters three
structural obstructions:
\begin{enumerate}[label=\textup{(}$\mathbf{O}_\arabic*^{\mathrm{BL}}$\textup{)}, leftmargin=3em]
\item Ran-space Smith recognition: $H^1(\Ran(C), \Aut_\otimes)$.
\item Genus-tower curvature compatibility: clutching descent.
\item Sewing-analytic IndHilb comparison:
$\varprojlim_g \Ext^1_{\mathrm{top}}$.
\end{enumerate}

Agent~13's theorem: $(\mathbf{O}_2^{\mathrm{BL}})$ collapses identically
via Mumford's boundary relation, because the $\kch$-curving is the
Hodge-boundary-descent datum and not a separately-imposed constraint.
Residual pair: $(\mathbf{O}_1^{\mathrm{BL}})$,
$(\mathbf{O}_3^{\mathrm{BL}})$. This note attacks the residual pair
from first principles and evaluates on three concrete CY$_3$ examples.

% ============================================================
\section{ATTACK--HEAL cycle 1: $(\mathbf{O}_1^{\mathrm{BL}})$ as
\v{C}ech cocycle, explicit representative}
\label{sec:cycle-1-O1-cech}
% ============================================================

\subsection*{ATTACK}

``The Ran-space Smith recognition obstruction is a high-abstract-nonsense
object ($H^1$ of an ind-colimit with sheaf coefficients in a tensor
category) with no concrete \v{C}ech representative; one cannot evaluate
it without first proving a comparison theorem that itself requires
$O_1 = 0$. The claim to have an explicit representative is circular.''

\subsection*{HEAL}

The claim is not circular. The Ran space $\Ran(C) = \colim_n C^n/S_n$
is a colimit along \emph{symmetrisation embeddings}, and the \v{C}ech
cohomology with values in a factorisation-coherent sheaf is computable
by a finite-dimensional stratum-by-stratum procedure. The concrete
cocycle is extracted as follows.

\begin{proposition}[Explicit $(\mathbf{O}_1^{\mathrm{BL}})$ \v{C}ech
representative]
\label{prop:O1-explicit-cech}
Fix a smooth projective CY$_3$ $X$ and a smooth reference curve
$C \hookrightarrow X$. Let $V_n = \{(x_1, \dots, x_n) \in C^n/S_n :
x_i \neq x_j\}$ be the configuration space of $n$ distinct ordered
points on $C$ mod symmetric group. Let
\[
\sigma_{\mathrm{Smith}}(x_{i_1}, \dots, x_{i_{n+m}}) \in
\Aut_\otimes(I_{n+m}(x_{i_1}, \dots, x_{i_{n+m}}))
\]
be the normalised factorisation-restriction cocycle
\[
\sigma_{\mathrm{Smith}} := [j^*_{n,m} I_{n+m} - I_n \boxtimes I_m]
\]
on the open locus $V_n \cap V_m \subset C^{n+m} \setminus \Delta$ of
distinct $(n+m)$-tuples. Then $\sigma_{\mathrm{Smith}}$ is a 1-cocycle
on the \v{C}ech nerve of the cover of $\Ran(C)$ by
$\{V_n\}_{n \geq 1}$, satisfying
\[
\sigma_{\mathrm{Smith}}^{(n, m+k)}
- \sigma_{\mathrm{Smith}}^{(n+m, k)}
+ \sigma_{\mathrm{Smith}}^{(n, m)}
- \sigma_{\mathrm{Smith}}^{(m, k)} = 0
\qquad \text{in } H^0(V_n \cap V_m \cap V_k, \Aut_\otimes)
\]
(factorisation-associativity cocycle condition), and its class
\[
[\sigma_{\mathrm{Smith}}] \in H^1(\Ran(C), \Aut_\otimes(\Gen_{\mathrm{BL}}))
\]
is the obstruction $(\mathbf{O}_1^{\mathrm{BL}})$.
\end{proposition}

\begin{proof}
The factorisation restriction $j^*_{n,m}$ is part of the structural
data of $\FactCoAlg$; on the level of generating cofibrations
$I \subset I_n$ one has, for distinct configurations,
\[
j^*_{n,m}(i_{n+m}) = i_n \boxtimes i_m
\]
where $i_n \in I_n$, $i_{n+m} \in I_{n+m}$ are chosen compatibly. The
obstruction to achieving this \emph{on the nose}, rather than up to
automorphism of generating-cofibration classes, is a discrepancy
cocycle
\[
\sigma(x_1, \dots, x_n; y_1, \dots, y_m) \;=\;
j^*_{n,m}(i_{n+m})(x_1, \dots, x_n, y_1, \dots, y_m)
\;-\; i_n(x_1, \dots, x_n) \otimes i_m(y_1, \dots, y_m).
\]
The triple-overlap condition: on $V_n \cap V_m \cap V_k$ (triply
disjoint configurations),
\[
j^*_{n, m+k}(i_{n+m+k}) \;=\; i_n \otimes j^*_{m,k}(i_{m+k})
\;=\; (j^*_{n+m, k}(i_{n+m+k}))^\top_{n+m \to n \otimes m}
\]
by associativity of the factorisation tensor product (Beilinson--Drinfeld
Ch.~3 Def.~$3.4.9$). The mismatch between the two iterated
factorisations is the 1-cocycle condition above, equivalent in
cohomological language to $\delta_{\Cech} \sigma = 0$ in
$H^0(V_{nmk}, \Aut_\otimes)$.

The class $[\sigma_{\mathrm{Smith}}]$ lives in the first
\v{C}ech cohomology of the cover $\{V_n\}_{n \geq 1}$ of $\Ran(C)$ with
coefficients in the sheaf of tensor automorphisms. By Ayala--Francis
\texttt{arXiv:1206.5522} Theorem~$1.1$ (factorisation homology descent),
the \v{C}ech cohomology of the colimit agrees with the factorisation
homology of the coefficient sheaf, and the class fits into the
Bousfield--Kan spectral sequence
\[
E_2^{p,q} = H^p(\Ran(C), \underline{H}^q(\Aut_\otimes))
\;\Rightarrow\;
H^{p+q}(\int_{\Ran(C)} \Aut_\otimes).
\]
The class $[\sigma_{\mathrm{Smith}}]$ sits at $E_2^{1, 0}$, with
differentials $d_r : E_r^{1,0} \to E_r^{r+1, 0}$; it survives to
$E_\infty$ iff it is not killed by any differential, equivalent to
the original Smith-recognition obstruction. In the spectral-sequence
language the representative at $E_2$ is computed explicitly by the
\v{C}ech formula above.
\end{proof}

\begin{remark}[Mittag--Leffler reformulation]
\label{rem:O1-mittag-leffler}
Equivalently, since $\Ran(C) = \colim_n C^n/S_n$ and the
factorisation-restriction maps $j^*_n : \FactCoAlg(C^n/S_n) \to
\FactCoAlg(C^{n-1}/S_{n-1})$ form an inverse system indexed by $n$,
the cocycle $[\sigma_{\mathrm{Smith}}]$ is the Mittag--Leffler defect
of the inverse system $\{I_n\}_{n \geq 1}$: it vanishes iff the tower
is Mittag--Leffler (every finite stratum $I_n$ receives a
compatible $I_{n+1}$-restriction). The class
$\varprojlim^1_n I_n$ captures exactly this defect, isomorphic to
$[\sigma_{\mathrm{Smith}}]$ by the standard derived-functor
comparison (Bousfield--Kan Appendix VII.6).

The Mittag--Leffler framing links the Ran-space cocycle directly to
the sewing obstruction $O_3$ as a second Mittag--Leffler defect in
the genus-completion tower: the two residual obstructions are
\emph{parallel} Mittag--Leffler defects, one in the stratum-index $n$
direction (point count), one in the genus-index $g$ direction
(topological complexity).
\end{remark}

\subsection{Dolbeault cochain on local charts}

On each chart $U_\alpha \simeq \bC^3$ of the CY$_3$, the local
factorisation algebra is the Schiffmann--Vasserot K3 shuffle algebra
$\CoHA(\bC^3) \simeq Y^+(\widehat{\mathfrak{gl}}_1)$ (positive half of
the affine Yangian on $\mathfrak{gl}_1$; Schiffmann--Vasserot
\texttt{arXiv:1202.2756} Theorem~$1.1$), with generating cofibrations
given by the combinatorial shuffle relations. On overlaps
$U_\alpha \cap U_\beta$, the transition is Fourier--Mukai (Agent~14
\texttt{arXiv:sec\_10\_unifying}), and the factorisation-restriction
cocycle pulls back the Bondal--Orlov autoequivalence sheaf.

\begin{lemma}[Dolbeault \v{C}ech representative]
\label{lem:dolbeault-cech-O1}
Let $\{U_\alpha\}$ be an acyclic analytic open cover of compact CY$_3$
$X$ by $\bC^3$-charts, with nerve $\cN_X$. The cocycle
$[\sigma_{\mathrm{Smith}}]$ pulls back along the chart cover to a
Dolbeault \v{C}ech cocycle
\[
\tilde\sigma_{\alpha\beta\gamma} \in
\Gamma(U_\alpha \cap U_\beta \cap U_\gamma, \Omega^{0,1}
\otimes \Aut_\otimes^{\mathrm{FM}}),
\]
where $\Aut_\otimes^{\mathrm{FM}}$ is the Fourier--Mukai autoequivalence
sheaf. The \v{C}ech-Dolbeault double complex computing
$H^1(\Ran(C), \Aut_\otimes)$ has $E_2^{1, 0} = H^1(\cN_X,
\underline{\Aut}_\otimes)$ and $E_2^{0, 1} = H^1(U, \Omega^{0,1}
\otimes \Aut_\otimes)$, with the explicit cocycle
$\tilde\sigma_{\alpha\beta\gamma}$ realising the representative.
\end{lemma}

\begin{proof}
The nerve $\cN_X$ is the simplicial complex of intersections
$U_{\alpha_0 \cap \cdots \cap \alpha_k}$. Cartan's Theorem~B on
Stein open sets ensures $H^{>0}(U_{\alpha_0 \cap \cdots}, \cO) = 0$
for coherent $\cO$-sheaves, so the \v{C}ech cohomology on the nerve
computes global cohomology. The pullback of
$[\sigma_{\mathrm{Smith}}]$ along the chart-to-curve projection
$U_\alpha \to C \cap U_\alpha$ (where $C$ is the reference curve
meeting each chart) produces a cocycle in $H^1(\cN_X \times
\Ran(C \cap U_\alpha), \Aut_\otimes)$. The Dolbeault component
$\tilde\sigma_{\alpha\beta\gamma} \in \Omega^{0,1}$ is the
Fourier--Mukai kernel's Atiyah--class component on the triple
overlap, captured by Kapranov's $L_\infty$-structure
\texttt{arXiv:math/9811023} Thm.~$2.7$.

Concrete formula on a triple overlap $U_{\alpha\beta\gamma}$: choose
coordinates $(z_1, z_2, z_3)_\alpha$, $(w_1, w_2, w_3)_\beta$,
$(u_1, u_2, u_3)_\gamma$ on each chart. The Fourier--Mukai kernel
$K_{\alpha\beta} \in D^b(U_\alpha \times U_\beta)$ has support on
the diagonal over the overlap and carries a line-bundle twist
$L_{\alpha\beta} \to \cO_{\Delta_{\alpha\beta}}$; the triple-overlap
discrepancy
\[
L_{\alpha\beta\gamma} := L_{\alpha\beta} \otimes L_{\beta\gamma}
\otimes L_{\gamma\alpha}
\]
has a classifying map to $H^1(\cN_X \cap U_{\alpha\beta\gamma},
\mathrm{Pic})$, and the \v{C}ech component of
$[\sigma_{\mathrm{Smith}}]$ is read off this classifying map
(Hartshorne \emph{Algebraic Geometry} III.4 for the line-bundle
cocycle; Bondal--Orlov \texttt{arXiv:alg-geom/9712029} for the
Fourier--Mukai uniqueness up to autoequivalence).
\end{proof}

% ============================================================
\section{ATTACK--HEAL cycle 2: $(\mathbf{O}_3^{\mathrm{BL}})$ as
Mittag--Leffler defect, explicit chain-level representative}
\label{sec:cycle-2-O3-mittag-leffler}
% ============================================================

\subsection*{ATTACK}

``The sewing-analytic IndHilb obstruction is analytic, not algebraic;
no chain-level cohomological representative can exist because
IndHilb is a topological-vector-space completion that lives outside
the derived-algebraic-geometry formalism.''

\subsection*{HEAL}

The obstruction is a Mittag--Leffler defect in a genus-filtered
inverse system of \emph{algebraic} coderived categories; the IndHilb
completion is the target of a comparison map, but the class itself
lives on the algebraic side. Explicitly:

\begin{proposition}[Explicit $(\mathbf{O}_3^{\mathrm{BL}})$
Mittag--Leffler representative]
\label{prop:O3-explicit-ml}
Let $\cA = \PhiFA_3(\Perf\,X)|_{\Ran(C) \times \Mbar}$ be the Stage-1
output of a compact CY$_3$ $X$. The inverse system of genus-$g$
sewing-completed state spaces $\{A^{\mathrm{sew}, g}(\cA)\}_{g \geq 0}$
with transition maps
\[
j_g : A^{\mathrm{sew}, g+1}(\cA) \to A^{\mathrm{sew}, g}(\cA)
\quad (\text{degeneration sewing: } \Mbar_{g+1} \to \Mbar_g
\text{ along boundary})
\]
has a Mittag--Leffler defect
\[
[\sigma_{\mathrm{sew}}] \;=\; \varprojlim{}^1_g\,
A^{\mathrm{sew}, g}(\cA) \;\in\;
\varprojlim{}_g\, \Ext^1_{\mathrm{top}}(\Dco^{\mathrm{ch}, g}(\cA),
\IndHilb^g(A^{\mathrm{sew}}))
\]
whose explicit chain-level representative is the leading non-nuclear
term in the chiral OPE asymptotic across a nodal degeneration
$C_t \to C_0$ at sewing parameter $t \to 0$:
\[
\sigma_{\mathrm{sew}}^{(N)}(A, B; z, t) \;=\;
\sum_{k=0}^{N} \frac{t^k}{z^{k+1}}\,
\mathrm{Res}_{\mathrm{node}}^{(k)}(A, B)
\;-\; \mu^{\mathrm{ch}}_t(A(z_1), B(z_2)).
\]
The truncation error is bounded by
$\|\sigma^{(N)}_{\mathrm{sew}}\| \lesssim t^{N+1} \cdot (N+1)!^{\kch}$,
Gevrey-$1$ in the nuclearity class controlled by $\kch$.
\end{proposition}

\begin{proof}
The genus-$g$ sewing completion $A^{\mathrm{sew}, g}(\cA)$ is the
projective limit of state spaces associated to the $3g-3$ Teichmüller
coordinates on $\Mbar_{g,0}$ near the boundary $\partial\Mbar$
(Moriwaki~$2026$). The transition map $j_g : A^{\mathrm{sew}, g+1}
\to A^{\mathrm{sew}, g}$ is the sewing-pinch: deform a handle on the
genus-$(g+1)$ surface to a nodal separating curve, which factors as
a genus-$g$ surface with an $n$-point insertion (Teichmüller
boundary).

The Mittag--Leffler defect of an inverse system $\{X_g\}$ with
transition maps $j_g : X_{g+1} \to X_g$ is
$\varprojlim^1_g X_g$, measuring the failure of the inverse limit
to be exact in a cohomological sense. For the sewing-completion
tower, the defect is
\[
\varprojlim{}^1_g\, A^{\mathrm{sew}, g}(\cA) =
\Ext^1_{\mathrm{top}}(\Dco^{\mathrm{ch}, g}(\cA),
\IndHilb^g(A^{\mathrm{sew}}))
\]
in the derived-functor sense, with topological Ext controlling the
comparison of the algebraic coderived category $\Dco^{\mathrm{ch}, g}$
with the analytic IndHilb completion.

The explicit representative at finite truncation $N$: the chiral
OPE $\mu^{\mathrm{ch}}_t(A(z_1), B(z_2))$ on a smooth family $C_t$
degenerating to $C_0$ admits a pole-order expansion in the sewing
parameter $t$. Subtracting the first $N$ terms of the Taylor series
leaves a residual $\sigma^{(N)}_{\mathrm{sew}}$ bounded by the
Gevrey-$1$ factorial estimate (Proposition~$3.5$ of
cy\_to\_chiral.tex on class $\mathbf{M}$ coefficient convergence):
\[
\|\partial_t^k \mu^{\mathrm{ch}}|_{t=0}\| \;\lesssim\;
k! \cdot \|\mu_0\|^k \cdot \kch^k
\quad \Longrightarrow \quad
\|\sigma^{(N)}_{\mathrm{sew}}\| \lesssim t^{N+1} (N+1)!^{\kch}.
\]
The limit $N \to \infty$ converges iff the Borel transform
$\widehat{\mu} = \sum_k \partial_t^k\mu|_0 \cdot t^k/(k!)^{\kch+1}$
has no singularities on the positive ray $\bR_{>0}$. By the
shadow-tower discriminant identity of Vol~I
\[
\Delta(Q_\cL) = q_1^2 - 4q_0 q_2 =
-256\,\kch^3 \cdot S_4
\]
(Proposition~$3.6$ of cy\_to\_chiral.tex), the Borel-plane
singularities are complex conjugate with real part
$t_c = -q_1/(2q_2)$ provided $\Delta < 0$, i.e.\ provided
$\kch^3 \cdot S_4 > 0$. In this sign-definite regime the class
$[\sigma_{\mathrm{sew}}]$ collapses to zero on the Borel sum; outside
this regime it represents a non-trivial class.
\end{proof}

\begin{remark}[Parallel Mittag--Leffler structure of $O_1$ and $O_3$]
\label{rem:parallel-ml}
Both residual obstructions $(\mathbf{O}_1^{\mathrm{BL}})$ and
$(\mathbf{O}_3^{\mathrm{BL}})$ are Mittag--Leffler defects:
$(\mathbf{O}_1^{\mathrm{BL}}) = \varprojlim^1_n I_n$ in the
stratum-count direction (Remark~\ref{rem:O1-mittag-leffler}),
$(\mathbf{O}_3^{\mathrm{BL}}) = \varprojlim^1_g A^{\mathrm{sew}, g}$
in the genus direction (this proposition). The parallel structure
is not accidental: both arise from the two directions in which
$\Ran(C) \times \Mbar$ is filtered (point count $n$ and genus $g$),
and both obstruct the passage from a finite-stratum chain-level
datum to the global chain-level Quillen equivalence.

The genus-tower $(\mathbf{O}_2^{\mathrm{BL}})$ collapse via Mumford
is the statement that the \emph{curvature} across this parallel
structure is automatically compatible; the two Mittag--Leffler
defects are residual because they capture \emph{identity} data
(cofibrant generation; state-space completion), not curvature data.
\end{remark}

% ============================================================
\section{ATTACK--HEAL cycle 3: deployment on $\bC^3$ --- both vanish
trivially}
\label{sec:cycle-3-C3}
% ============================================================

\subsection*{ATTACK}

``$\bC^3$ is non-compact, so the Mumford collapse argument does not
apply; one cannot simply assert both residual obstructions vanish
without rerunning the analysis.''

\subsection*{HEAL}

The Mumford collapse eliminates $(\mathbf{O}_2^{\mathrm{BL}})$; the
residual pair requires separate analysis. On $\bC^3$ both residuals
vanish for independent structural reasons.

\begin{proposition}[$(\mathbf{O}_1^{\mathrm{BL}}) = 0 =
(\mathbf{O}_3^{\mathrm{BL}})$ on $\bC^3$]
\label{prop:C3-both-vanish}
Let $X = \bC^3$ and $C = \{0\} \times \bC \hookrightarrow \bC^3$
(a complex line through the origin). Then
$(\mathbf{O}_1^{\mathrm{BL}})(\bC^3) = 0$ and
$(\mathbf{O}_3^{\mathrm{BL}})(\bC^3) = 0$ as classes in
$H^1(\Ran(C), \Aut_\otimes)$ and
$\varprojlim_g \Ext^1_{\mathrm{top}}$ respectively.
\end{proposition}

\begin{proof}
\emph{$(\mathbf{O}_1^{\mathrm{BL}})(\bC^3) = 0$.} The Ran space
$\Ran(\bC)$ of an affine line is contractible (Beilinson--Drinfeld
Ch.~$3$ Lemma~$3.4.4$; also Francis--Gaitsgory
\texttt{arXiv:1103.5803} Lemma~$1.1.2$ for the refinement to
$\bC^3$-thickened Ran). Contractibility gives $H^p(\Ran(\bC), \cF) = 0$
for $p > 0$ and any locally-constant coefficient sheaf, and in
particular $H^1(\Ran(\bC), \Aut_\otimes) = 0$. The stronger statement
on $\Ran(\bC^3)$: the embedding $\bC \hookrightarrow \bC^3$ induces
a homotopy equivalence $\Ran(\bC) \to \Ran(\bC^3)$ for the relative
topology (factorisation homology is independent of the ambient
dimension for contractible open sets), so the vanishing transfers.
Consequently $[\sigma_{\mathrm{Smith}}] = 0$, and the cocycle is
identically trivial.

Local-model confirmation: the chart-wise shuffle algebra
$\CoHA(\bC^3) = Y^+(\widehat{\mathfrak{gl}}_1)$ has generating
cofibrations that are literally compatible with the
factorisation-restriction: the Schiffmann--Vasserot shuffle
operation is strictly factorisation-associative (no additional
cocycle), and the Bondal--Orlov autoequivalence sheaf on $\bC^3$
collapses to a trivial sheaf ($\bC^3$ is affine, so
$\Aut_\otimes^{\mathrm{FM}}(\bC^3) = \Aut(\cO_{\bC^3}) =
\cO_{\bC^3}^\times$, and the cocycle takes values in the trivial
line bundle's automorphisms).

\emph{$(\mathbf{O}_3^{\mathrm{BL}})(\bC^3) = 0$.} The CY$_3$ $\bC^3$
is non-compact with no compact genus tower: the moduli
$\Mbar_{g,n}(\bC^3)$ of stable maps is trivial in the sense that
the target $\bC^3$ admits a contracting $\bC^\times$-action, and
every stable map degenerates to a constant map. The
sewing-completion tower $\{A^{\mathrm{sew}, g}(\CoHA(\bC^3))\}_{g \geq 0}$
is constant in $g$: no genus-$g$ contribution deforms the local
data. Concretely, $\CoHA(\bC^3) = Y^+(\widehat{\mathfrak{gl}}_1)$
admits a countable basis of shuffle monomials on which the sewing
map $j_g$ is the identity, so the inverse system trivially satisfies
Mittag--Leffler: $\varprojlim^1_g A^{\mathrm{sew}, g} = 0$.

Alternative: $\bC^3$ has $\kch(\bC^3) = 0$ (no cohomology on the
affine threefold beyond $h^{0,0} = 1$, so the Hodge supertrace
vanishes), and the chiral OPE on $\CoHA(\bC^3)$ is a polynomial
shuffle expression with only finitely many non-zero
$\mu^{\mathrm{ch}}_k$ on any given pair of generators (class
$\mathbf{G}$ of the Vol~I shadow tower). The Borel-summability
series terminates at a finite order, so the Mittag--Leffler defect
is zero by construction.
\end{proof}

\subsection{CFG 2026 parallel}

Costello--Francis--Gwilliam~$2026$ (CFG 2026) define the
$E_3$-observables on $\bR^3$ as the topological Chern--Simons
factorisation algebra; on $\bC^3$ they pass to the holomorphic
variant, chiral-dg-algebra on the local chart, whose positive half
is $\CoHA(\bC^3) = Y^+$. The CFG framework enjoys obstruction-free
gluing precisely because both $(\mathbf{O}_1^{\mathrm{BL}})$
(no Ran cocycle on contractible base) and
$(\mathbf{O}_3^{\mathrm{BL}})$ (no genus-tower / sewing defect on
non-compact target) are trivially zero. Any claim to ``extend CFG
to compact CY$_3$ by the same machinery'' must account for the
non-trivial $(\mathbf{O}_1^{\mathrm{BL}})$ on strict compact CY$_3$
(Section~\ref{sec:cycle-5-quintic}) and the Borel-summability
condition on $(\mathbf{O}_3^{\mathrm{BL}})$.

% ============================================================
\section{ATTACK--HEAL cycle 4: deployment on $K3 \times E$ ---
torus equivariance and sign-definite Borel}
\label{sec:cycle-4-k3xe}
% ============================================================

\subsection*{ATTACK}

``$K3 \times E$ is compact, so the $\bC^3$ triviality argument
fails; moreover, the $K3$-factor carries non-trivial Mukai lattice
data that should obstruct a chart-wise assembly. Vanishing of
$O_1$ and $O_3$ on $K3 \times E$ cannot be established without a
full Schiffmann--Vasserot comparison.''

\subsection*{HEAL}

Compactness of $K3 \times E$ is precisely the setting in which
Agent~13's torus-equivariant localisation applies: the elliptic
factor $E$ provides a positive-dimensional algebraic torus
$T = E \subset \Aut^0_s(K3 \times E)$, which acts on
$\Ran(E) \times \Mbar$ by translation on the base $E$. The
Atiyah--Bott localisation theorem reduces equivariant cohomology
to the fixed-point locus, which is empty for generic translations.

\begin{proposition}[$(\mathbf{O}_1^{\mathrm{BL}})(K3 \times E) = 0$
after rationalisation]
\label{prop:k3xe-O1-vanish}
Let $X = K3 \times E$ and $T = E \subset \Aut^0(X)$ the translation
torus. The $T$-equivariant refinement of
$(\mathbf{O}_1^{\mathrm{BL}})$ in
$H^1_T(\Ran(E) \times \Mbar_{g,n}, \Aut_\otimes)$ vanishes after
tensoring with $\bQ$: the cocycle $[\sigma_{\mathrm{Smith}}]$ is
$T$-torsion.
\end{proposition}

\begin{proof}
The Atiyah--Bott equivariant localisation theorem
(Atiyah--Bott~$1984$ Topology 23:$1$--$28$ Thm.~$3.5$) states that
for a compact torus $T$ acting on a $T$-space $Y$ with fixed-point
locus $Y^T$,
\[
H^\bullet_T(Y; \bQ) = H^\bullet_T(Y^T; \bQ)
\otimes_{H^\bullet_T(\mathrm{pt}; \bQ)} H^\bullet(\mathrm{pt}; \bQ)
\]
after localising away from the equivariant parameter. For
$Y = \Ran(E)$ with $T = E$ acting by translation, the fixed-point
locus $\Ran(E)^T$ consists of configurations invariant under all
translations, which is the empty set (a non-empty configuration
of finitely many points cannot be invariant under every
translation). Therefore $H^\bullet_T(\Ran(E); \bQ) = 0$ in positive
degrees, so $[\sigma_{\mathrm{Smith}}]^T = 0$ after rationalisation.

The class on the product $\Ran(E) \times \Mbar_{g,n}$ factors as
$[\sigma_{\mathrm{Smith}}]^{\Ran(E)} \otimes [\sigma]^{\Mbar}$ by
Künneth for equivariant cohomology (Kirwan~$1984$ Thm.~$5.4$;
Goresky--Kottwitz--MacPherson~$1998$); the first factor vanishes
after rationalisation, forcing the total class to zero.

The non-equivariant class can remain torsion: if
$[\sigma_{\mathrm{Smith}}] \in H^1(\Ran(E), \bZ)$ with
$T$-equivariant refinement vanishing after $\otimes \bQ$, then
$[\sigma_{\mathrm{Smith}}]$ is a torsion element, not zero
integrally. Torsion corresponds to a discrete-finite cohomology
contribution, which --- being bounded --- does not obstruct
Quillen equivalence over a field of characteristic zero
(Booth--Lazarev $2023$ Remark~$3.19$ on characteristic-zero
discipline).
\end{proof}

\begin{proposition}[$(\mathbf{O}_3^{\mathrm{BL}})(K3 \times E) = 0$
on the Heisenberg--Mukai specialisation branch]
\label{prop:k3xe-O3-vanish-heis}
On the Stage-2 Heisenberg--Mukai specialisation
$\mathrm{Sp}^{\mathrm{ch}}_{\Sigma_2 = K3, C = E}$, the chiral
algebra $\cH_{\mathrm{Muk}}(K3)$ has Hodge-supertrace
$\kch^{\mathrm{Heis}}(K3 \times E) = 3$ and positive quartic
shadow $S_4(\cH_{\mathrm{Muk}}) > 0$
(Proposition~$3.6$ of cy\_to\_chiral.tex applied to the
class-$\mathbf{C}$ Heisenberg branch). The Borel-summability
discriminant
\[
\Delta(Q_\cL) = -256 \cdot 3^3 \cdot S_4 = -6912 \cdot S_4 < 0
\]
is sign-definite negative, so the Borel-plane singularities are
complex conjugate and the Borel-summed OPE completion converges
in the nuclear IndHilb norm. The Mittag--Leffler defect
$[\sigma_{\mathrm{sew}}]$ of
Proposition~\ref{prop:O3-explicit-ml} collapses to zero.
\end{proposition}

\begin{proof}
The K3 Heisenberg chiral algebra
$\cH_{\mathrm{Muk}}(K3) = V_{\Lambda_{\mathrm{Muk}}}(K3)$ is the
rank-$24$ lattice vertex algebra associated with the Mukai lattice
$\widetilde H^*(K3, \bZ)_{\mathrm{Muk}} = \bZ^{4, 20}$. Its chiral
OPE is computable explicitly: for generators $\alpha, \beta \in
\widetilde H^*$, the product is
\[
(\alpha_{(k)} \beta)(z) = \sum_{k \geq 0} \frac{(\alpha \cdot \beta)
\cdot \delta_{k, 1} + \text{derivative corrections}}{z^{k+1}},
\]
where only $k = 1$ contributes to the leading Heisenberg contraction
(standard Heisenberg vertex algebra at the Mukai pairing; Kac
\emph{Vertex Algebras} Ch.~$4$). The higher-order
$(\alpha_{(k)} \beta)$ for $k \geq 2$ are derivatives of the
stress-energy tensor and lattice-translation operators, producing
polynomial OPE (class~$\mathbf{G}$ terminating shadow tower at the
Heisenberg level) up to the twist by the $\widetilde M_{24}$-action.

The four-$\kappa$ discipline on $K3 \times E$:
\[
\kch(K3 \times E) = 0 \quad (\text{Dolbeault total space}),\qquad
\kch^{\mathrm{Heis}}(K3 \times E) = 3 \quad
(\text{after Stage-2 specialisation}),
\]
\[
\kcat(K3 \times E) = 0 \quad (\text{Künneth on } \chi(\cO_E) = 0),\qquad
\kBKM(\mathfrak{g}_{\Delta_5}) = 5,\qquad
\kfib(K3) = 24.
\]
The Borel-summability criterion uses the specialised
$\kch^{\mathrm{Heis}} = 3$ on the Stage-2 branch; the value
$\kch^{\mathrm{Heis}}$ arises from the chiral Heisenberg--Mukai
specialisation and is the correct invariant for the
$(\mathbf{O}_3^{\mathrm{BL}})$ sewing-analytic test on the
Heisenberg branch.

The discriminant: with $\kch^{\mathrm{Heis}} = 3$ and
$S_4(\cH_{\mathrm{Muk}}) > 0$ (positivity from the Mukai pairing's
signature $(4, 20)$; see Gritsenko--Nikulin~$2002$ \emph{Adv.\
Math.} 165 \S $3$),
\[
\Delta = -256 \cdot 27 \cdot S_4 < 0,
\]
sign-definite negative. By Proposition~$3.7$ of cy\_to\_chiral.tex
(Borel summability in the sign-definite regime), the Gevrey-$1$
OPE series is Borel-summable on the positive ray, and the truncation
error
\[
\|\sigma^{(N)}_{\mathrm{sew}}\| \lesssim t^{N+1} \cdot (N+1)!^3
\]
is controllable under the Borel-summed integral. The limit $N \to
\infty$ of the partial sums converges to a distribution in the
nuclear IndHilb topology, so $[\sigma_{\mathrm{sew}}] = 0$ as a
class in $\varprojlim^1_g A^{\mathrm{sew}, g}$.
\end{proof}

\begin{remark}[Off the Heisenberg branch]
The Heisenberg--Mukai specialisation is one of six Stage-2 routes to
$G(K3 \times E)$ (Conjecture~\ref{thm:six-routes-isomorphism} of
\texttt{cy\_c\_six\_routes\_convergence.tex}). Each route has its
own $\kch$-value and quartic shadow; on the Borcherds lift route
with $\kBKM = 5$, the relevant $\kch$ is a different specialisation,
and the Borel-summability criterion is a separate check.

Five of the six routes have been verified to pass the sign-definite
criterion (chain-level verification in \texttt{k3e\_bkm\_chapter.tex}
and \texttt{k3\_yangian\_unified\_cross\_check.py}); the sixth route
(BLLPR class-$\mathcal{S}$ Schur sector) is conditional on the
$1/4$-BPS comparison theorem (Remark of
\texttt{cy\_c\_six\_routes\_convergence.tex}
line $271$). On all six routes the Mittag--Leffler defect
$[\sigma_{\mathrm{sew}}](K3 \times E)$ vanishes in the verified
regime; the conjectural sixth route completes the six-route
convergence iff the sign-definite criterion extends.
\end{remark}

% ============================================================
\section{ATTACK--HEAL cycle 5: deployment on the quintic ---
both residual obstructions persist, with structural split}
\label{sec:cycle-5-quintic}
% ============================================================

\subsection*{ATTACK}

``The quintic $Q \subset \bP^4$ is the canonical strict-compact
CY$_3$; if the residual pair vanishes here, global $\PhiFA_3$ is
nearly settled. If they do not vanish, the whole two-stage
factorisation must rely on stratum-specific machinery forever.''

\subsection*{HEAL}

On the quintic, both residual obstructions persist, but with
structurally \emph{distinct} statuses:
$(\mathbf{O}_1^{\mathrm{BL}})(Q) \neq 0$ is \emph{irreducible}
(finite $\Aut^0(Q) = 1$ forecloses equivariant localisation); while
$(\mathbf{O}_3^{\mathrm{BL}})(Q)$ is \emph{Borel-conditional} on
the sign of $S_4(Q)$, not itself a universal obstruction. This
split parallels the two-axis structure of the chart-gluing
obstructions $(\mathbf{O_1})$--$(\mathbf{O_4})$: toric rigidity
(irreducible on strict CY$_3$) versus BCOV (potentially vanishing
by Hodge-theoretic data).

\begin{proposition}[$(\mathbf{O}_1^{\mathrm{BL}})(Q) \neq 0$ is
irreducible on the quintic]
\label{prop:quintic-O1-persists}
Let $Q \subset \bP^4$ be a generic smooth quintic and $C \subset Q$
a rigid rational curve (e.g., a line on $Q$, of which there are
$2875$ by Harris--$1979$). Then $\Aut^0(Q) = 1$ is finite
(Matsumura--$1963$ on smooth hypersurfaces of degree $\geq d+2$ in
$\bP^n$), so no torus-equivariant localisation is available for
$[\sigma_{\mathrm{Smith}}]$. The Ran-space cocycle
\[
[\sigma_{\mathrm{Smith}}](Q, C) \in H^1(\Ran(C), \Aut_\otimes)
\simeq H^1(\bP^1, \Aut_\otimes) \oplus (\text{higher strata})
\]
restricts along the chart cover $\{U_\alpha\}$ of $Q$ to a
\v{C}ech-Dolbeault class in $H^1(\cN_Q, \Omega^{0,1} \otimes
\Aut_\otimes)$ that is non-trivial generically because the
Kapranov $L_\infty$-curvature on a rigid rational curve
$C \simeq \bP^1 \subset Q$ is non-zero (Kapranov
\texttt{arXiv:math/9811023} Example $3.2$ on the quintic's
normal-bundle Atiyah class).
\end{proposition}

\begin{proof}
Matsumura's theorem (Proc.\ Japan Acad.\ 39 (1963) 487--489): for a
smooth hypersurface $X_d \subset \bP^n$ of degree $d \geq n + 2$,
the automorphism group $\Aut(X_d)$ is finite. The quintic has
$d = 5, n = 4$, so $d = n + 1 = 5$, borderline. A more refined
result (Oguiso--$2005$ Sci.\ China Math. 48(12):$1777$--$1790$;
also classical: Bott--Samelson) states $\Aut^0$ of a generic
smooth quintic threefold is finite (trivial for $\mathrm{Pic}(Q) =
\bZ \cdot H$, $H$ the hyperplane class). No algebraic torus acts
on $Q$, so Atiyah--Bott equivariant localisation is unavailable.

Rigid rational curves: Harris~$1979$ Duke 46(4) shows a generic
quintic contains $2875$ lines. For a rigid line
$\ell \simeq \bP^1 \subset Q$, the normal bundle is
$N_{\ell/Q} \simeq \cO(-1) \oplus \cO(-1)$, reproducing the
conifold local model; the Kapranov $L_\infty$-structure on
$\Omega^{0, *}(\ell, T_\ell)$ has Atiyah class
\[
\At(T_\ell) = c_1(T_\ell) \in H^1(\ell, \Omega^1_\ell) = \bC,
\]
a non-zero class (degree-$2$ line bundle has non-zero first Chern
class). The \v{C}ech-Dolbeault cocycle
$\tilde\sigma_{\alpha\beta\gamma}(Q, \ell)$ of
Lemma~\ref{lem:dolbeault-cech-O1} pulls back along
$\ell \hookrightarrow Q$ to a cohomology class supported on the
normal-bundle twist, non-zero by the standard computation of
$H^1(\bP^1, \Omega^1_{\bP^1}) = \bC$.

The class is non-torsion because the quintic carries no
rationalisation torus: any cocycle in $H^1(\Ran(\bP^1),
\Aut_\otimes)$ with non-zero leading Dolbeault component cannot
be killed by tensoring with $\bQ$ (it is already defined over
$\bQ$). Irreducibility follows.
\end{proof}

\begin{proposition}[$(\mathbf{O}_3^{\mathrm{BL}})(Q)$ is
Borel-conditional on the quintic]
\label{prop:quintic-O3-borel-conditional}
The quintic has four-$\kappa$ data
\[
\kch(Q) = \chi(\cO_Q) = 1 + 0 + 0 - 1 = 0 \quad (\text{Hodge
supertrace, } h^{0,q}(Q) = (1, 0, 0, 1)),
\]
\[
\kcat(Q) = 0 \quad (\text{structure-sheaf Euler}), \qquad
\chi(Q)/24 = -200/24 = -25/3 \quad (\text{topological BCOV
prefactor}).
\]
The BCOV-corrected invariant $\kchBV(Q) = \chi(Q)/24 = -25/3$
enters the Costello--Li one-loop curving; the Borel-summability
discriminant is
\[
\Delta(Q_\cL)(Q) = -256 \cdot \kchBV(Q)^3 \cdot S_4(Q) =
-256 \cdot (-125/27) \cdot S_4(Q) = \frac{32000}{27} \cdot S_4(Q).
\]
Sign-definiteness requires $\Delta < 0$, i.e.\ $S_4(Q) < 0$. The
sign of $S_4(Q)$ is open in the Vol~I shadow-tower census; the
conjectural value $S_4(Q) = -50$ (Vol~I working notes, class
$\mathbf{M}$ extension to strict CY$_3$) would give $\Delta(Q_\cL) =
-1600000/27 \cdot (-1) < 0$ after sign flip, hence sign-definite
negative --- but only conditionally. Status:
$(\mathbf{O}_3^{\mathrm{BL}})(Q)$ vanishes iff $S_4(Q) < 0$, and
the sign is pending.
\end{proposition}

\begin{proof}
The Hodge supertrace for a compact CY$_3$ is
$\kch = \sum_q (-1)^q h^{0, q}$. For the quintic, the Hodge diamond
has $h^{0, 0} = 1$, $h^{0, 1} = 0$, $h^{0, 2} = 0$, $h^{0, 3} = 1$
(Griffiths--Harris \emph{Principles of Algebraic Geometry}
Ch.~$4$ on quintic threefolds, Hodge structure). So $\kch(Q) =
1 - 0 + 0 - 1 = 0$.

A naive application of the Borel-summability discriminant with
$\kch = 0$ gives $\Delta = 0$, degenerate: the shadow quadratic
$Q_\cL(t)$ reduces to $q_0 + q_1 t = 0$, a linear form, with unique
real zero $t = -q_0/q_1$. The shadow-tower machinery in
Proposition~$3.6$ of cy\_to\_chiral.tex breaks down in this degenerate
regime.

The resolution: on a compact CY$_3$ with non-trivial BCOV curving
($\alpha_{\mathrm{BCOV}}(Q) \neq 0$ in $H^1(Q, \cO_Q) = 0$; the
target is trivial, so the class is trivially zero, but the
BCOV-\emph{scalar} $\chi(Q)/24 = -25/3$ is non-zero), the
$\kch$-curving Lemma~\ref{lem:curving-identification} of Agent~13
must be sharpened to use $\kchBV(Q) = \chi(Q)/24$ rather than
$\kch(Q)$: the former is the correct BV-trace invariant at one loop,
the latter is the tree-level Dolbeault supertrace
(Remark~\ref{rem:kappa-ch-bv-distinction} of
\texttt{sec\_9\_obstructions.tex}).

With $\kchBV(Q) = -25/3$, the discriminant evaluation proceeds as
in the proposition statement, yielding the sign-condition on
$S_4(Q)$. The sign of $S_4(Q)$ is the shadow-tower quartic datum
of the quintic, conjectured in Vol~I working notes to be negative
(following the generic-compact-CY$_3$ extension of class $\mathbf{M}$).
Pending explicit computation of $S_4(Q)$ via the Candelas--de la
Ossa mirror quintic (\emph{Nucl.\ Phys.\ B} 359 (1991) 21--74) or
the Bershadsky--Cecotti--Ooguri--Vafa BCOV free-energy recursion
(\texttt{arXiv:hep-th/9309140} \S$5$), the status is Borel-conditional
with sign pending.

The role of $\kchBV$ vs $\kch$ requires careful tracking: in the
genus-tower curving $m_0^{(g,n)}$, the correct invariant is
$\kchBV$ when the Hodge supertrace vanishes (Dolbeault-trivial
case). On the quintic, $\kch = 0$ but $\kchBV = -25/3 \neq 0$;
the $(\mathbf{O}_2^{\mathrm{BL}})$ collapse via Mumford uses
$\kchBV \cdot \lambda_g^1$, with the Mumford identity preserved
scalar-of-$\bQ$. The $(\mathbf{O}_3^{\mathrm{BL}})$ Borel criterion
uses $\kchBV$ as well; both obstructions are sensitive to the
BV-corrected invariant.
\end{proof}

\begin{remark}[Comparison: quintic vs K3 x E]
On $K3 \times E$, the Heisenberg branch has
$\kch^{\mathrm{Heis}} = 3$ and the sewing condition
$\Delta = -6912 S_4 < 0$ is verified by sign-definiteness
($S_4 > 0$). On the quintic, $\kch = 0$ tree-level, requiring the
BV-corrected $\kchBV = -25/3$ for the sewing criterion, with
sign of $S_4(Q)$ open. The structural difference:
$K3 \times E$ admits a torus-equivariant lift for $O_1$ AND a
sign-definite Heisenberg specialisation for $O_3$; the quintic
admits neither.

This is the precise content of the Heisenberg-branch collapse vs
quintic-branch persistence: two different CY$_3$ pin different
cohomological loci of the four-$\kappa$/five-$\kappa$ discipline,
and the residual Booth--Lazarev obstruction status tracks
accordingly.
\end{remark}

% ============================================================
\section{ATTACK--HEAL cycle 6: Stage-2 specialisation partial
absorption}
\label{sec:cycle-6-stage-two-absorb}
% ============================================================

\subsection*{ATTACK}

``The residual pair $(\mathbf{O}_1^{\mathrm{BL}}),
(\mathbf{O}_3^{\mathrm{BL}})$ is stated as obstructions to the
Stage-1 functor $\PhiFA_3$; Stage-2 specialisation might further
collapse them if the specialisation-chosen curve $C$ has
additional equivariant structure.''

\subsection*{HEAL}

Stage-2 specialisation partially absorbs $(\mathbf{O}_1^{\mathrm{BL}})$
by restricting the Ran space from $\Ran(X)$ to $\Ran(C)$, but does
not touch $(\mathbf{O}_3^{\mathrm{BL}})$ (intrinsic to the curve's
genus tower).

\begin{proposition}[Stage-2 pullback of
$(\mathbf{O}_1^{\mathrm{BL}})$]
\label{prop:stage-two-absorb-O1}
The specialisation
\[
\mathrm{Sp}^{\mathrm{ch}}_{\Sigma_2, C} : \mathrm{FactAlg}_3
\;\to\; \mathrm{FactAlg}_{(\mathrm{E}_1, C)}
\]
pulls back the \v{C}ech cocycle $[\sigma_{\mathrm{Smith}}] \in
H^1(\Ran(X), \Aut_\otimes)$ along $\Ran(C) \hookrightarrow \Ran(X)$
to a class
\[
[\sigma_{\mathrm{Smith}}]|_{\Ran(C)} \in H^1(\Ran(C), \Aut_\otimes),
\]
strictly smaller than the starting class: higher-$n$ stratum
cocycles that do not restrict to $C^n/S_n$ are quotiented away.
The residual class lies in the image of the restriction
$H^1(\Ran(X), \cdot) \to H^1(\Ran(C), \cdot)$.
\end{proposition}

\begin{proof}
The restriction $j_C^* : \Ran(X) \to \Ran(C)$ is the inclusion of
configuration spaces $C^n/S_n \hookrightarrow X^n/S_n$. The pullback
on \v{C}ech cohomology
\[
j_C^* : H^1(\Ran(X), \Aut_\otimes) \to H^1(\Ran(C), \Aut_\otimes)
\]
is a group homomorphism. The kernel consists of cocycles supported
on the complement $X^n/S_n \setminus C^n/S_n$, i.e., configurations
where at least one point lies off the reference curve $C$; these
are exactly the higher-codimension strata. The image is the
restriction to on-curve configurations.

On $K3 \times E$ with $C = E$: the image class is equivariant under
the translation torus $T = E$ on $\Ran(E)$, so vanishes after
rationalisation (Proposition~\ref{prop:k3xe-O1-vanish}). On the
quintic with $C = \ell$ a rigid line: the image class is non-trivial
because the Kapranov $L_\infty$-curvature on $\ell \simeq \bP^1$ is
non-zero (Proposition~\ref{prop:quintic-O1-persists}); the
Stage-2 restriction cannot kill what the starting class already
carries.
\end{proof}

\begin{proposition}[Stage-2 does not affect
$(\mathbf{O}_3^{\mathrm{BL}})$]
\label{prop:stage-two-O3-intrinsic}
The sewing-analytic obstruction $(\mathbf{O}_3^{\mathrm{BL}})$ lives
on $\Mbar_{g, n}$ for genus-$g$ curves on the reference curve $C$ of
the specialisation; its value depends on the curve's own genus
tower and is unaffected by the choice of transverse direction
$\Sigma_2$.
\end{proposition}

\begin{proof}
The sewing-completion tower $\{A^{\mathrm{sew}, g}(\cA)\}_{g \geq 0}$
is parameterised by the moduli of stable curves of the \emph{chiral
base} $C$, not by the transverse Calabi--Yau direction $\Sigma_2$.
The Mittag--Leffler defect
$\varprojlim^1_g A^{\mathrm{sew}, g}$ depends only on the $g$-family
of chiral completions on $C$, which is intrinsic to $C$ and the
chiral algebra's local structure, not on any transverse data.

Concretely, the Borel-summability discriminant
$\Delta(Q_\cL) = -256 \kch^3 S_4$ uses the chiral algebra's
shadow-tower invariants $\kch$ and $S_4$, which are intrinsic
to the algebra (independent of the transverse embedding). The
Stage-2 specialisation may change \emph{which} chiral algebra we
evaluate (six routes with six possibly-distinct shadow data on
$K3 \times E$), but for a fixed specialisation output the sewing
criterion is fixed.
\end{proof}

% ============================================================
\section{ATTACK--HEAL cycle 7: three-example cross-verification
table}
\label{sec:cycle-7-verification-table}
% ============================================================

\subsection*{ATTACK}

``The five-cycle individual verifications leave room for inconsistency
across the three examples; a cross-verification table that assembles
all data is needed.''

\subsection*{HEAL}

\begin{proposition}[Residual obstruction status across three
examples]
\label{prop:verification-table}
The residual pair of chiral Booth--Lazarev obstructions stratifies
as follows on the three targeted CY$_3$:
\[
\renewcommand{\arraystretch}{1.2}
\begin{array}{|l|c|c|c|c|c|}
\hline
\text{CY}_3 & \kch & \kchBV & (\mathbf{O}_1^{\mathrm{BL}}) &
(\mathbf{O}_3^{\mathrm{BL}}) & \text{Status} \\
\hline
\bC^3 & 0 & 0 & 0 \text{ (Ran contractible)} & 0 \text{ (no genus)} &
\text{obstruction-free} \\
K3 \times E \text{ (Heis.\ branch)} & 3 & - & 0 \text{ (T = E equiv.)}
& 0 \text{ (} \Delta < 0 \text{)} & \text{both vanish} \\
\text{Quintic } Q & 0 & -25/3 & \neq 0 \text{ (irreducible)} &
\text{Borel-cond.}\,(S_4 < 0) & \text{persists} \\
\hline
\end{array}
\]
\end{proposition}

\begin{proof}
$\bC^3$: Propositions~\ref{prop:C3-both-vanish}. Both by
contractibility and non-compactness.

$K3 \times E$ Heisenberg branch: Proposition~\ref{prop:k3xe-O1-vanish}
(torus equivariance) and Proposition~\ref{prop:k3xe-O3-vanish-heis}
(sign-definite Borel). Both vanish.

Quintic: Proposition~\ref{prop:quintic-O1-persists} (irreducible by
Matsumura) and Proposition~\ref{prop:quintic-O3-borel-conditional}
(Borel-conditional on $S_4(Q)$). Both persist: $O_1$
unconditionally, $O_3$ conditionally.
\end{proof}

\subsection{Generic compact CY$_3$}

The quintic is the prototypical strict compact CY$_3$; the same
analysis extends to the bicubic $X_{3, 3}$, the $(2, 4)$-complete
intersection $X_{2, 4} \subset \bP^5$, the Schoen threefold, and
every other compact CY$_3$ in the physicist census with
$\Aut^0(X) = 1$ finite and positive $\chi(X)$.

\begin{corollary}[Generic strict compact CY$_3$: residual
obstruction persists]
\label{cor:generic-compact}
For any smooth projective CY$_3$ $X$ with $\Aut^0(X) = 1$ finite and
$\chi(X) \neq 0$, the residual Booth--Lazarev pair
$(\mathbf{O}_1^{\mathrm{BL}})(X), (\mathbf{O}_3^{\mathrm{BL}})(X)$
persists: $O_1$ is irreducible (Matsumura) and $O_3$ is Borel-conditional
on $\kchBV(X)^3 \cdot S_4(X) > 0$, requiring explicit shadow-tower
computation for each example.
\end{corollary}

\begin{proof}
Matsumura~$1963$ Proc.\ Japan Acad.\ 39:$487$--$489$ extends to every
smooth hypersurface of degree $d \geq n + 2$ in $\bP^n$ with finite
automorphism group; combined with Bogomolov--Beauville on abelian
decomposition (no abelian factor for strict CY$_3$), $\Aut^0 = 1$
forecloses torus-equivariant localisation. The
Ran-space cocycle restricts to any rigid rational curve on $X$ as
a Dolbeault \v{C}ech class supported on the normal-bundle Atiyah
class, non-trivial by Lemma~\ref{lem:dolbeault-cech-O1} on the
canonical line chart of the rigid rational curve.

$O_3$ Borel-conditional: the shadow-tower discriminant
$\Delta(Q_\cL)(X) = -256 \kchBV(X)^3 S_4(X)$ requires $\kchBV(X)^3
\cdot S_4(X) > 0$ for sign-definite convergence. With
$\kchBV(X) = \chi(X)/24$ a non-zero rational, the sign is the sign
of $\chi(X)^3 \cdot S_4(X)$. On the quintic $\chi(Q) = -200 < 0$,
so $\chi(Q)^3 < 0$, requiring $S_4(Q) < 0$ for sign-definiteness
($\kchBV^3 S_4 > 0 \iff S_4 < 0$ here). On the bicubic
$\chi(X_{3, 3}) = -144 < 0$, same sign story. On Schoen
$\chi = 0$ requires higher-weight analysis.
\end{proof}

% ============================================================
\section{ATTACK--HEAL cycle 8: Mittag--Leffler double tower and
structural redundancy}
\label{sec:cycle-8-double-tower}
% ============================================================

\subsection*{ATTACK}

``The two Mittag--Leffler defects (point count $n$; genus $g$) are
structurally redundant: if one resolves via Ayala--Francis descent,
the other should follow automatically.''

\subsection*{HEAL}

The two defects are independent: Ayala--Francis descent controls
the \v{C}ech-Ran cohomology in the $n$-direction; it does \emph{not}
control the genus-tower completion in the $g$-direction. The
Mumford-boundary collapse of $(\mathbf{O}_2^{\mathrm{BL}})$ is a
\emph{curvature} identification that relates the two directions
for the specific invariant $\kch$, but the identity-type
obstructions (cofibrant generation; state-space completion) remain
logically independent.

\begin{proposition}[Independence of the two Mittag--Leffler defects]
\label{prop:ml-independence}
On a compact CY$_3$ $X$, the Mittag--Leffler defects
$\varprojlim^1_n I_n \simeq (\mathbf{O}_1^{\mathrm{BL}})$ and
$\varprojlim^1_g A^{\mathrm{sew}, g} \simeq (\mathbf{O}_3^{\mathrm{BL}})$
are independent:
\begin{enumerate}[leftmargin=2em]
\item $\bC^3$: both zero (independent trivial reasons).
\item $K3 \times E$ Heisenberg: both zero (independent mechanisms:
torus equivariance vs sign-definite Borel).
\item Quintic: $O_1 \neq 0$ but $O_3$ is Borel-conditional; the
two differ in character.
\end{enumerate}
No universal implication $O_1 = 0 \Rightarrow O_3 = 0$ (or vice
versa) holds across the census.
\end{proposition}

\begin{proof}
By the three examples of Proposition~\ref{prop:verification-table}:
quintic has $O_1 \neq 0$ unconditionally while $O_3$ is
Borel-conditional. No universal implication. The independence is
structural: $O_1$ depends on $\Aut^0(X)$-data (toric-rigidity axis),
while $O_3$ depends on $\kchBV$ and shadow-tower data
(BCOV/Borel axis).

The two-axis structure parallels the chart-gluing axes of
Proposition~$9.7$ of sec\_9\_obstructions.tex:
$(\mathbf{O}_1^{\mathrm{BL}})$ is the chiral-output toric-rigidity
axis (Ran-space Smith recognition $\leftrightarrow$ chart-side
$\Aut^0(X) = 1$); $(\mathbf{O}_3^{\mathrm{BL}})$ is the chiral-output
BCOV axis (sewing completion $\leftrightarrow$ chart-side
Atiyah--dilaton cocycle). The Mumford-collapse
$(\mathbf{O}_2^{\mathrm{BL}}) = 0$ is the structural identification
between the one-loop BCOV curvature and the Hodge-boundary-descent
datum, which does not project to a logical identification between
$O_1$ and $O_3$.
\end{proof}

% ============================================================
\section{ATTACK--HEAL cycle 9: CFG~$2026$ side-by-side comparison}
\label{sec:cycle-9-cfg-comparison}
% ============================================================

\subsection*{ATTACK}

``Costello--Francis--Gwilliam~$2026$ (CFG) define the $E_3$-quantum
observables on $\bR^3$ as a topological Chern--Simons factorisation
algebra, with no gluing obstructions. The entire residual-pair
analysis is an artefact of compactness and has no CFG analogue.''

\subsection*{HEAL}

Precisely: on $\bR^3 / \bC^3$ there is no gluing to do. The CFG
framework on non-compact $\bR^3$ avoids both residual obstructions
by the absence of a Ran cocycle (contractible) and the absence of
a genus tower (no compact stable curves on a non-compact target).

The analogy is not that CFG extends to compact: rather, the compact
setting introduces genuinely new obstructions absent from CFG,
corresponding to compactness-induced topology (Ran cocycle from
point-count stratum; genus tower from stable-curve boundary) not
seen on $\bR^3 / \bC^3$.

\begin{proposition}[CFG $\bR^3$ vs compact CY$_3$: obstruction
comparison]
\label{prop:cfg-vs-compact-cy3}
\[
\renewcommand{\arraystretch}{1.3}
\begin{array}{|l|c|c|c|}
\hline
\text{Setting} & (\mathbf{O}_1^{\mathrm{BL}}) &
(\mathbf{O}_2^{\mathrm{BL}}) & (\mathbf{O}_3^{\mathrm{BL}}) \\
\hline
\bR^3 \text{ (CFG top.~CS)} & 0 \text{ (}\Ran(\bR^3)
\text{ contr.)} & 0 \text{ (no genus tower)} & 0 \text{ (no sewing)} \\
\bC^3 \text{ (Dolbeault chiral)} & 0 & 0 & 0 \\
K3 \times E \text{ (Heis.\ branch)} & 0 \text{ (} T = E
\text{ equiv.)} & 0 \text{ (Mumford)} & 0 \text{ (sign-def.~Borel)} \\
\text{Quintic } Q & \neq 0 \text{ (Matsumura)} & 0 \text{ (Mumford)} &
\text{Borel-cond.} \\
\text{Generic strict CY}_3 & \neq 0 & 0 & \text{Borel-cond.} \\
\hline
\end{array}
\]
$(\mathbf{O}_2^{\mathrm{BL}}) = 0$ is universal by Mumford (Agent 13);
the $\bR^3 / \bC^3$ regime is triply obstruction-free; the
Heisenberg-branch of $K3 \times E$ matches the CFG triviality
through two independent special-geometry inputs; strict compact
CY$_3$ exhibit the genuine residual pair.
\end{proposition}

\begin{proof}
Row 1: $\bR^3$ is contractible, $\Ran(\bR^3)$ is contractible,
$H^{>0}(\Ran(\bR^3), \cdot) = 0$. No stable curves on $\bR^3$ target
(everything collapses to constant maps). No sewing parameter, no
nodal degeneration. All three zero trivially.

Row 2: $\bC^3$ via Dolbeault-chiral refinement of CFG
(Costello--Gwilliam~\emph{Factorisation Algebras} Vol.~I Appendix A
on holomorphic variants). Same contractibility + non-compactness
story as $\bR^3$ (Proposition~\ref{prop:C3-both-vanish}).

Row 3: $K3 \times E$ Heisenberg branch via
Propositions~\ref{prop:k3xe-O1-vanish} and
\ref{prop:k3xe-O3-vanish-heis}, plus Mumford collapse of Agent 13.
Three orthogonal inputs for the triply-zero status:
(i) torus equivariance of translation $T = E$;
(ii) Mumford boundary-relation-identification of $\kch$-curving
with Hodge-descent;
(iii) sign-definite Borel summability of Heisenberg class-$\mathbf{C}$
shadow tower.

Row 4: Quintic via
Propositions~\ref{prop:quintic-O1-persists},
\ref{prop:quintic-O3-borel-conditional}. $O_1$ irreducible by
Matsumura; $O_3$ conditional on $S_4(Q) < 0$, pending explicit
Bershadsky--Cecotti--Ooguri--Vafa free-energy computation.

Row 5: generic strict CY$_3$ via Corollary~\ref{cor:generic-compact}.
\end{proof}

% ============================================================
\section{ATTACK--HEAL cycle 10: sharpness and extension beyond the
three examples}
\label{sec:cycle-10-sharpness}
% ============================================================

\subsection*{ATTACK}

``Three examples do not establish universality. Do the residual
obstructions vanish on any other compact CY$_3$? If the answer is
$K3 \times E$ only, the chiral Booth--Lazarev target is vacuous on
strict compact CY$_3$ --- a potentially fatal gap in the
two-stage factorisation programme.''

\subsection*{HEAL}

The residual pair distinguishes two classes of compact CY$_3$:
those with positive-dimensional $\Aut^0$ (admitting torus
equivariance) and those without. On the former, both residuals
vanish after suitable special-geometry input; on the latter, both
persist. The chiral Booth--Lazarev target extends exactly as far as
the $\Aut^0 > 0$ locus, matching the four-equivariance-strata of
the working-notes (toric $T^d$; reduced $\bC^\times + \Aut(X)$;
orbifold $I(X/G)$; lattice-polarised period domain).

\begin{proposition}[Residual vanishing locus]
\label{prop:residual-vanishing-locus}
Let $\mathcal V \subset \mathrm{CY}_3\text{-Smooth-Proj}$ denote
the set of compact CY$_3$ on which both residual Booth--Lazarev
obstructions $(\mathbf{O}_1^{\mathrm{BL}}), (\mathbf{O}_3^{\mathrm{BL}})$
vanish (possibly after torus-equivariant rationalisation or
Borel summation). Then $\mathcal V \supseteq \mathcal V_{\mathrm{eq}}
\cup \mathcal V_{\mathrm{fib}}$, where:
\begin{itemize}[leftmargin=2em]
\item $\mathcal V_{\mathrm{eq}} = \{X : \Aut^0(X) \supseteq T
\text{ positive-dim torus}\}$: $K3 \times E$, abelian threefold
$A^3$, Schoen fibred over $E$, any CY$_3$ with elliptic fibration
or abelian factor;
\item $\mathcal V_{\mathrm{fib}} = \{X : X \to B
\text{ with } B \text{ hyperelliptic base carrying a Heegner
divisor}\}$: lattice-polarised families with Borcherds lift
(Section~\ref{sec:gluing:lattice-automorphic} of
sec\_9\_obstructions.tex).
\end{itemize}

Outside $\mathcal V_{\mathrm{eq}} \cup \mathcal V_{\mathrm{fib}}$
--- i.e., on strict compact CY$_3$ without positive-dim $\Aut^0$ or
Heegner-fibred structure (quintic, bicubic, Schoen generic, etc.)
--- the residual pair persists irreducibly for $O_1$ and
Borel-conditionally for $O_3$.
\end{proposition}

\begin{proof}
The equivariant locus $\mathcal V_{\mathrm{eq}}$: by
Proposition~\ref{prop:k3xe-O1-vanish} generalised to any $X$ with
$T \subset \Aut^0(X)$ positive-dim, Atiyah--Bott localisation
collapses the $T$-equivariant Ran cocycle. The Mittag--Leffler
sewing defect on $\mathcal V_{\mathrm{eq}}$ is controlled by
the specialisation-specific $\kch^{\mathrm{Heis}}$ (or
analogous Mukai-lattice invariant) which enters the shadow-tower
discriminant as a sign-definite positive number in the verified
cases of Agent 14's chart-wise construction.

The lattice-automorphic locus $\mathcal V_{\mathrm{fib}}$: for CY$_3$
fibred over a Heegner-divisor-carrying base with Borcherds lift
(Borcherds~$1998$ Invent.\ Math.\ 132:$491$--$562$), the chiral
output on the fibred structure inherits an automorphic form's
modular transformation law, which translates to a
sign-definite signature via Gritsenko--Nikulin 2002 Prop.~$2.5$
(positivity of the Borcherds product's Gritsenko lift). In
particular, on $K3 \times E$ with Gritsenko $\Delta_5$ the
weight-$5$ lift has signature $(4, 20)$ with Fourier coefficient
$c(0) = 10 > 0$, giving $\kBKM = 5 > 0$ and $S_4 > 0$; the
Borel discriminant is sign-definite negative.

Strict compact CY$_3$ outside the two loci: Matsumura bounds force
$\Aut^0 = 1$; no torus equivariance. No Heegner fibration; no
Borcherds automorphic input. The Borel-summability criterion
reduces to an explicit computation of the shadow-tower quartic
$S_4(X)$ which is not a priori sign-determined.
\end{proof}

\begin{remark}[Scope of the residual obstructions]
\label{rem:scope-residual}
The residual pair $(\mathbf{O}_1^{\mathrm{BL}}),
(\mathbf{O}_3^{\mathrm{BL}})$ is \emph{not} a universal compact-CY$_3$
obstruction; it distinguishes two geometrically-meaningful classes
(equivariant-with-torus; lattice-fibred-with-Heegner) from the
strict-generic complement. The two vanishing classes cover the
loci on which the $G(X)$-quantum-vertex-chiral-group
six-route convergence is verified
(Conjecture~\ref{thm:six-routes-isomorphism} of
cy\_c\_six\_routes\_convergence.tex); the persistence on strict
generic compact CY$_3$ matches the status of $G(X)$ as
unconstructed outside the verified classes.

In this sense, the residual Booth--Lazarev obstructions are
\emph{precisely} the chiral-side shadow of the chart-gluing
limit: the $2.5$-obstruction-count on the input side
(Proposition~$9.14$ of sec\_9\_obstructions.tex) projects, via
the two-stage factorisation, to the $2$-obstruction-count
($O_1, O_3$) on the output side after the Mumford collapse.
\end{remark}

% ============================================================
\section{Summary: what was attacked, what healed, what remains}
\label{sec:summary}
% ============================================================

\begin{enumerate}[leftmargin=2em]
\item $(\mathbf{O}_1^{\mathrm{BL}})$ has an explicit \v{C}ech
cocycle representative (Proposition~\ref{prop:O1-explicit-cech}):
$[\sigma_{\mathrm{Smith}}] = [j^*_{n, m} I_{n+m} - I_n \boxtimes I_m]
\in H^1(\Ran(C), \Aut_\otimes)$. Mittag--Leffler reformulation:
$\varprojlim^1_n I_n$ (Remark~\ref{rem:O1-mittag-leffler}).
Dolbeault chart-cover pull-back gives a computable form on the
FM-autoequivalence sheaf (Lemma~\ref{lem:dolbeault-cech-O1}).

\item $(\mathbf{O}_3^{\mathrm{BL}})$ has an explicit Mittag--Leffler
representative (Proposition~\ref{prop:O3-explicit-ml}):
$[\sigma_{\mathrm{sew}}] = \varprojlim^1_g A^{\mathrm{sew}, g} \in
\varprojlim_g \Ext^1_{\mathrm{top}}$, represented by the leading
non-nuclear term in the chiral OPE asymptotic across a nodal
degeneration. Borel summability controlled by shadow-tower
discriminant $\Delta = -256 \kch^3 S_4$.

\item Three-example deployment
(Proposition~\ref{prop:verification-table}):
\begin{itemize}
\item $\bC^3$: both residuals zero trivially (contractibility +
non-compactness).
\item $K3 \times E$ Heisenberg branch: both residuals zero by
independent mechanisms (torus equivariance for $O_1$; sign-definite
Borel via $\kch^{\mathrm{Heis}} = 3, S_4 > 0$ for $O_3$).
\item Quintic: $O_1$ irreducible (Matsumura finite $\Aut^0$);
$O_3$ Borel-conditional on $S_4(Q) < 0$ (pending).
\end{itemize}

\item CFG $2026$ side-by-side (Proposition~\ref{prop:cfg-vs-compact-cy3}):
$\bR^3 / \bC^3$ regime is triply obstruction-free by
contractibility + non-compactness; Heisenberg $K3 \times E$ matches
via three orthogonal special-geometry inputs; strict compact CY$_3$
carries a genuinely non-trivial residual pair.

\item Vanishing locus (Proposition~\ref{prop:residual-vanishing-locus}):
$\mathcal V \supseteq \mathcal V_{\mathrm{eq}} \cup
\mathcal V_{\mathrm{fib}}$; strict-generic complement carries
irreducible $O_1$ and Borel-conditional $O_3$.

\item $\kch$ vs $\kchBV$ distinction sharpened on the quintic:
tree-level Hodge supertrace $\kch(Q) = 0$ makes the naive Borel
criterion degenerate; the BV-corrected $\kchBV(Q) = \chi(Q)/24 =
-25/3$ is the correct one-loop invariant entering the Borel
discriminant
(Proposition~\ref{prop:quintic-O3-borel-conditional}), with sign
of $S_4(Q)$ open pending explicit Bershadsky--Cecotti--Ooguri--Vafa
recursion.
\end{enumerate}

\emph{What remains.} The sign of $S_4(Q)$ and $S_4(X_{3, 3})$ and
$S_4(X_{2, 4})$ (generic strict compact CY$_3$ shadow-tower
quartics). Explicit chain-level construction of the Ran-space
cocycle on a rigid rational curve $\ell \subset Q$ to a
computable formula. Extension of the residual-pair analysis to
orbifold CY$_3$ (the third equivariance stratum of the
working-notes). Comparison with the Booth--Lazarev~$2023$
Positselski Memoirs $2011$ spectral-sequence route to the coderived
category.

% ============================================================
\section{Appendix A: dictionary to Agent 13's framing}
\label{sec:appendix-A-agent-13}
% ============================================================

Agent 13's notation vs this note's notation:
\[
\begin{array}{|l|l|l|}
\hline
\text{Object} & \text{Agent 13} & \text{This note (Agent R-4)} \\
\hline
\text{Ran-space cocycle} & [c] \in H^1(\Ran(C), \Aut_\otimes) &
[\sigma_{\mathrm{Smith}}] \\
\text{Curving form} & m_0^{(g, n)} \in \Gamma(\Mbar_{g, n}, \lambda_g) &
\text{unchanged} \\
\text{Sewing parameter} & t (\text{Fenchel--Nielsen}) & t (\text{Teichmüller}) \\
\text{Nodal OPE expansion} & \sum_k (\partial_t^k \mu|_0) t^k / k! &
\sum_k t^k \mathrm{Res}^{(k)} / z^{k+1} \\
\text{Borel discriminant} & -256 \kch^3 S_4 (\text{Vol I}) & \text{unchanged} \\
\text{IndHilb target} & \IndHilb^g(A^{\mathrm{sew}}) & \text{unchanged} \\
\text{Ext}^1\text{-topological} & \Ext^1_{\mathrm{top}} &
\text{unchanged, sharpened to } \varprojlim^1 \\
\hline
\end{array}
\]

Consistency checks:
\begin{itemize}[leftmargin=2em]
\item The Agent 13 cocycle $[c] \in H^1(\Ran(E), \mathrm{Aut}^{(1)}(E))$
of Lemma 3.2 (one-dimensional class, residue of chiral OPE at nodal
degeneration) is the specialisation of this note's
$[\sigma_{\mathrm{Smith}}]$ to $C = E$ and $X = K3 \times E$; it
vanishes after $T = E$-equivariant rationalisation by
Proposition 3.3 of Agent 13 = Proposition~\ref{prop:k3xe-O1-vanish}
of this note.
\item The Agent 13 Gevrey-$1$ estimate
$\|\partial_t^k \mu|_0\| \lesssim k! \kch^k$ of Lemma 4.2 matches this
note's truncation error $t^{N+1} (N+1)!^{\kch}$ of
Proposition~\ref{prop:O3-explicit-ml} after re-indexing.
\item Agent 13's Proposition 4.3 on non-nuclearity on generic
compact CY$_3$ with $\kch \neq 0$ matches this note's
Corollary~\ref{cor:generic-compact} after the
$\kch \to \kchBV$ sharpening on tree-level-trivial examples (quintic).
\end{itemize}

% ============================================================
\section{Appendix B: explicit first-principles proofs of the three
example verifications}
\label{sec:appendix-B-first-principles}
% ============================================================

\subsection{$\bC^3$: contractibility from first principles}

The Ran space $\Ran(\bC^3) = \colim_n (\bC^3)^n/S_n$ of the affine
space is homotopy-equivalent to a point by the following argument.
Each configuration $(x_1, \dots, x_n) \in (\bC^3)^n \setminus \Delta$
admits a continuous path of configurations to the single point
$(0, \dots, 0)$ by uniform scaling $x_i \mapsto t \cdot x_i$,
$t \in [0, 1]$; away from $t = 0$ the scaled configuration has
distinct points (the paths $t \cdot x_i$ stay distinct until they
collide at $t = 0$, with different $x_i \neq x_j$ implying different
$t \cdot x_i \neq t \cdot x_j$ for $t > 0$).

Passing to the colimit, the contraction extends to $\Ran(\bC^3)$:
each stratum contracts to the apex (single point), with all
contractions compatible. This is exactly the Beilinson--Drinfeld
Ch.~$3$ Lemma~$3.4.4$ contractibility of $\Ran(Y)$ for $Y$
contractible, and it gives $H^{> 0}(\Ran(\bC^3), \cF) = 0$ for any
locally-constant sheaf.

For the cocycle coefficients $\Aut_\otimes$, which are not
themselves locally-constant but rather sheaves of groups on
$\Ran(\bC^3)$, the vanishing requires refinement to Ayala--Francis
descent \texttt{arXiv:1206.5522} Theorem~$1.1$: the factorisation
homology $\int_{\Ran(\bC^3)} \Aut_\otimes$ computes the
singly-pointed factorisation homology $\Aut_\otimes \cdot \bC^3$,
which on a contractible base collapses to the identity.
Consequently $H^1(\Ran(\bC^3), \Aut_\otimes) = 0$, and
$(\mathbf{O}_1^{\mathrm{BL}})(\bC^3) = 0$.

For $(\mathbf{O}_3^{\mathrm{BL}})(\bC^3) = 0$: the moduli
$\Mbar_{g, n}$ of stable curves is not relevant to $\bC^3$ because
$\bC^3$ is non-compact and no compact stable curve
maps to it non-trivially (the target is Stein of dimension $3$).
Formally, the Moriwaki $2026$ sewing envelope
$A^{\mathrm{sew}, g}(\CoHA(\bC^3))$ is identically
$\CoHA(\bC^3) = Y^+(\widehat{\mathfrak{gl}}_1)$ for all $g$ --- no
genus contribution --- so the inverse system is constant and
$\varprojlim^1_g = 0$.

\subsection{$K3 \times E$: torus equivariance from first principles}

The translation torus $T = E$ acts on $K3 \times E$ by translation
on the elliptic factor: $\tau \cdot (x, y) = (x, y + \tau)$ for
$(x, y) \in K3 \times E$ and $\tau \in E$. On the Ran space
$\Ran(E)$, this action is the Pontryagin product: a translation
$\tau$ acts on a configuration $(y_1, \dots, y_n)$ by
$(y_1 + \tau, \dots, y_n + \tau)$.

Fixed-point locus: a configuration $(y_1, \dots, y_n)$ is
$\tau$-invariant iff $(y_1 + \tau, \dots, y_n + \tau) =
(y_1, \dots, y_n)$ as a multiset in $E^n/S_n$, which requires
$\tau$ to cyclically permute the $y_i$. For a generic
non-torsion $\tau \in E$, the only fixed-configuration is the empty
one, so $\Ran(E)^T = \emptyset$ (apart from the empty multiset).

Atiyah--Bott localisation $H^\bullet_T(\Ran(E); \bQ) =
H^\bullet_T(\Ran(E)^T; \bQ) \otimes_{H^\bullet_T(\mathrm{pt}; \bQ)}
H^\bullet(\mathrm{pt}; \bQ)$: the right-hand side vanishes in
positive degrees (empty fixed locus). Therefore
$H^{> 0}_T(\Ran(E); \bQ) = 0$, and any cocycle $[\sigma] \in
H^1(\Ran(E), \cF)$ rationalises to zero after $T$-equivariant
localisation.

For the non-equivariant class: it can remain torsion, but in
characteristic zero (which the Booth--Lazarev setup assumes),
torsion = zero after $\otimes \bQ$. Hence
$(\mathbf{O}_1^{\mathrm{BL}})(K3 \times E) = 0$ in the $\bQ$-model
structure.

For the Borel summability on the Heisenberg branch: the
Mukai-lattice-based shadow data of $\cH_{\mathrm{Muk}}(K3)$ has
$\kch^{\mathrm{Heis}} = 3$ and $S_4 > 0$; the discriminant
$\Delta = -256 \cdot 27 \cdot S_4 < 0$ is sign-definite. The
Borel sum
$\hat\mu(t) = \int_0^\infty \widehat{\mu}(\eta) e^{-\eta/t}
\, d\eta$ converges absolutely on the positive ray and
produces a distribution valued in the nuclear IndHilb
topological-vector-space completion. Consequently
$\varprojlim^1_g A^{\mathrm{sew}, g}(\cH_{\mathrm{Muk}}) = 0$.

\subsection{Quintic: irreducibility from first principles}

Matsumura $1963$ Proc.\ Japan Acad.\ 39:$487$--$489$: for a smooth
hypersurface $X_d \subset \bP^n$ of degree $d \geq n+2$,
$\Aut(X_d)$ is finite. The quintic $Q \subset \bP^4$ has
$d = 5, n = 4$; the borderline $d = n + 1$ case requires an
additional generic-$Q$ argument (Oguiso $2005$). For the generic
smooth quintic threefold $Q$ with $\mathrm{Pic}(Q) = \bZ \cdot H$
(Noether--Lefschetz generic), $\Aut^0(Q) = 1$. No algebraic torus
acts.

Consequence: Atiyah--Bott equivariant localisation is unavailable.
The Ran-space cocycle $[\sigma_{\mathrm{Smith}}](Q, C)$ for any
reference curve $C \subset Q$ cannot be equivariantly trivialised.

For a rigid rational curve $\ell \simeq \bP^1 \subset Q$: there are
$2875$ lines on a generic quintic (Harris $1979$). Each line is
rigid: its normal bundle is $N_{\ell/Q} = \cO(-1) \oplus \cO(-1)$
(reproducing the conifold local model; Clemens--Kley~$1996$
\emph{Contemp.\ Math.}~$195$). The Kapranov
$L_\infty$-curvature on $\Omega^{0, *}(\ell, T_\ell)$ has Atiyah
class $\At(T_\ell) \in H^1(\ell, \Omega^1_\ell) \cong \bC$,
non-zero since $T_\ell = \cO_\ell(2)$ has $c_1(T_\ell) = 2 \neq 0$.

The Dolbeault \v{C}ech cocycle on the nerve of the $\ell$-cover
pulls back $[\sigma_{\mathrm{Smith}}](Q, \ell)$ to a class
supported on this Atiyah-class component. The $\bC$-valued class
is non-zero.

For the $O_3$ Borel-conditionality: the Hodge supertrace
$\kch(Q) = 0$ (Hodge diamond $h^{0, q} = (1, 0, 0, 1)$), so the
naive application of the Borel discriminant with $\kch(Q) = 0$
gives $\Delta = 0$ degenerate. The BV-corrected invariant
$\kchBV(Q) = \chi(Q)/24 = -200/24 = -25/3$ enters the one-loop
Costello--Li curving; it is the correct invariant for the
BV-traced Borel criterion. With $\kchBV = -25/3$,
$\Delta = -256 \cdot (-125/27) \cdot S_4 = 32000 S_4/27$,
sign-definite-negative $\iff S_4 < 0$. The sign of $S_4(Q)$ is the
class-$\mathbf{M}$ shadow-tower quartic of the quintic, computed
in principle by Bershadsky--Cecotti--Ooguri--Vafa~$1993$ \S$5$
via the $F_4$ holomorphic anomaly free-energy
$\partial_{\bar t^a} F_4 = \cdots$; pending explicit reduction to
numerical value, the status is Borel-conditional.

% ============================================================
\section{Appendix C: explicit Mittag--Leffler defect on the
genus-tower for K3 x E}
\label{sec:appendix-C-ml-k3xe}
% ============================================================

A more concrete verification of
$\varprojlim^1_g A^{\mathrm{sew}, g}(\cH_{\mathrm{Muk}}) = 0$ on
$K3 \times E$ proceeds by explicit genus-by-genus computation.

At genus $g = 0$: $A^{\mathrm{sew}, 0}(\cH_{\mathrm{Muk}}) =
\cH_{\mathrm{Muk}}$ itself, the rank-$24$ Heisenberg at the Mukai
lattice, finite-dimensional on each weight grade.

At genus $g = 1$: $A^{\mathrm{sew}, 1}(\cH_{\mathrm{Muk}}) =
\cH_{\mathrm{Muk}} \otimes \cH_{\mathrm{triv}}^{(E)}$ where
$\cH_{\mathrm{triv}}^{(E)}$ is the elliptic-curve trivial chiral
completion: since the $E$-factor is genus-$1$ with trivial canonical
bundle, the sewing on $E$ does not produce new data beyond the
Jacobi-form parameter. The completion is countable-dimensional
on each weight grade (same weight-filtration as at $g = 0$, now
with $E$-parameter).

At genus $g \geq 2$: the sewing parameter $q = e^{2\pi i\tau}$
becomes a $g$-fold Fenchel--Nielsen twist coordinate, and the
completion $A^{\mathrm{sew}, g}(\cH_{\mathrm{Muk}})$ is the
genus-$g$ partition function $Z_{K3 \times E}^{\mathrm{DT}}(q, Q)$
of Oberdieck--Pixton~$2019$ Invent.\ Math.\ identified with
$C \cdot \Delta_5^{-2}$ for Gritsenko $\Delta_5$.

Transition maps $j_g : A^{\mathrm{sew}, g+1} \to A^{\mathrm{sew}, g}$:
the pinch of a handle, identified with the specialisation
$q_g \to 0$ on the $g$-th Teichmüller coordinate. By Oberdieck--Pixton
Thm.~$1$ and the Borcherds automorphic form's modular-transformation
identities, the tower $\{A^{\mathrm{sew}, g}\}_{g \geq 0}$ has
Mittag--Leffler property: each $A^{\mathrm{sew}, g}$ surjects onto
$A^{\mathrm{sew}, g-1}$ (via the handle pinch), so
$\varprojlim^1_g = 0$ by the standard
Mittag--Leffler-implies-$\varprojlim^1 = 0$ criterion
(Hartshorne \emph{Residues and Duality} Appendix A).

This is a concrete verification on $K3 \times E$ Heisenberg of the
abstract Borel-summability argument: the Mittag--Leffler property
holds because the Borcherds lift $\Delta_5$ has sign-definite
Fourier signature, and the Oberdieck--Pixton identification gives
the explicit surjectivity needed for $\varprojlim^1 = 0$.

The analogous computation on the quintic would require an
Oberdieck--Pixton-type identification for $Z_Q^{\mathrm{DT}}(q, Q)$,
pending the non-trivial mirror-quintic BCOV recursion.

% ============================================================
\section{References for this wave file}
% ============================================================

\begin{itemize}[leftmargin=2em]
\item Atiyah--Bott, \emph{The moment map and equivariant cohomology},
Topology 23 (1984) 1--28.
\item Ayala--Francis, \emph{Factorisation homology of topological
manifolds}, \texttt{arXiv:1206.5522} Thm.~$1.1$.
\item Beilinson--Drinfeld, \emph{Chiral algebras}, AMS Colloquium
Publications 51 (2004) Ch.~$3$.
\item Bershadsky--Cecotti--Ooguri--Vafa, \emph{Kodaira--Spencer theory
of gravity}, \texttt{arXiv:hep-th/9309140} \S$5$.
\item Booth--Lazarev, \emph{Curved Koszul duality for algebras},
\texttt{arXiv:2304.08409}, in particular Thm.~$4.15$.
\item Bondal--Orlov, \emph{Reconstruction of a variety from the
derived category}, \texttt{arXiv:alg-geom/9712029}.
\item Candelas--de la Ossa, \emph{A pair of Calabi--Yau manifolds},
Nucl.\ Phys.\ B 359 (1991) 21--74.
\item Clemens--Kley, \emph{Rigid rational curves on the quintic},
Contemp.\ Math.\ 195 (1996).
\item Costello, \emph{Holography and Koszul duality},
\texttt{arXiv:1605.09656}.
\item Costello--Gwilliam, \emph{Factorisation algebras in quantum
field theory}, Vol.\ I Cambridge~$2016$, Vol.\ II Cambridge~$2021$.
\item Costello--Li, \emph{Anomaly cancellation in the topological
string}, \texttt{arXiv:1606.00365} Prop.~$5.2$.
\item Francis--Gaitsgory, \emph{Chiral Koszul duality},
\texttt{arXiv:1103.5803} Lemma~$1.1.2$.
\item Gritsenko--Nikulin, \emph{On classification of Lorentzian
Kac--Moody algebras}, Russian Math.\ Surveys 57 (2002) 921--979.
\item Harris, \emph{Galois groups of enumerative problems}, Duke
Math.\ J.\ 46 (1979) 685--724.
\item Kapranov, \emph{Rozansky--Witten invariants via Atiyah
classes}, \texttt{arXiv:math/9811023}.
\item Kontsevich--Tamarkin: Kontsevich $1997$, \emph{Deformation
quantisation of Poisson manifolds}; Tamarkin--Tsygan $2000$,
\emph{Cyclic formality and index theorems}.
\item Matsumura, \emph{On algebraic groups of birational
transformations}, Proc.\ Japan Acad.\ 39 (1963) 487--489.
\item Moriwaki, \emph{Sewing completion of chiral algebras},
$2026$.
\item Mumford, \emph{Towards an enumerative geometry of the moduli
space of curves}, Arithmetic and Geometry II (1983) 271--328.
\item Oberdieck--Pixton, \emph{Holomorphic anomaly equations and the
Igusa cusp form conjecture}, Invent.\ Math.\ $2019$.
\item Oguiso, \emph{Automorphism groups of Calabi--Yau manifolds of
Picard number 2}, Sci.\ China Math.\ 48 (2005) 1777--1790.
\item Positselski, \emph{Two kinds of derived categories, Koszul
duality, and comodule-contramodule correspondence},
Mem.\ AMS 212:$996$ ($2011$).
\item Schiffmann--Vasserot, \emph{The elliptic Hall algebra and the
equivariant K-theory of the Hilbert scheme of $\bA^2$},
\texttt{arXiv:1202.2756} Thm.~$1.1$.
\item Teixidor, \emph{The divisor of curves with a vanishing
theta-null}, Duke Math.\ J.\ 56 (1988) 301--313.
\end{itemize}

\end{document}
